\chapter{Proporcionalidade}\label{ch:g7:prop}

O Ano 6 conheceu as tabelas \omterm{def:g6:propdata:def}{proporcionais} (\cref{ch:g6:propdata});
este capítulo transforma a \omterm{def:g7:prop:table}{proporcionalidade} numa ferramenta completa:
a regra da cruz para encontrar um quarto valor, as porcentagens e as
\omterm{def:g7:prop:scale}{escalas} de mapas e maquetes. A velocidade, a mais famosa de todas as
\omterm{def:g7:prop:table}{proporcionalidades}, é estudada no \cref{ch:g8:speed}.

\section{Reconhecer e completar tabelas}

\begin{definition}[Tabela de proporcionalidade]\label{def:g7:prop:table}
Uma tabela de duas linhas é uma \emph{tabela de
proporcionalidade}\index{proporcionalidade} quando os números da
segunda linha são os da primeira multiplicados por um número fixo, o
\emph{coeficiente}. Dito de outro modo: todos os \omterm{def:g3:division:remainder}{quocientes} de coluna
$\frac{\text{de baixo}}{\text{de cima}}$ são iguais.
\end{definition}

\begin{example}\label{ex:g7:prop:recognize}
\begin{center}
\renewcommand{\arraystretch}{1.3}
\begin{tabular}{l|ccc}
de cima & $4$ & $10$ & $14$ \\
\hline
de baixo & $6$ & $15$ & $21$ \\
\end{tabular}
\end{center}
\omterm{def:g3:division:remainder}{Quocientes}: $\frac64 = 1.5$, $\frac{15}{10} = 1.5$,
$\frac{21}{14} = 1.5$: proporcional, com coeficiente $1.5$.
\end{example}

\begin{theorem}[Regra da cruz]\label{thm:g7:prop:cross}
Numa \omterm{def:g7:prop:table}{tabela de proporcionalidade} com colunas $\begin{pmatrix} a \\ b
\end{pmatrix}$ e $\begin{pmatrix} c \\ d \end{pmatrix}$, os \omterm{def:g2:mult:def}{produtos}
cruzados são iguais:
\[
a \times d = b \times c .
\]
Por consequência, o quarto valor pode ser calculado a partir dos outros
três: $d = \dfrac{b \times c}{a}$.
\end{theorem}

\begin{proof}
Seja $k$ o coeficiente: $b = ka$ e $d = kc$. Então
\[
a \times d = a \times kc = k \times ac,
\qquad
b \times c = ka \times c = k \times ac :
\]
os dois \omterm{def:g2:mult:def}{produtos} cruzados valem $k \times ac$. Dividindo $ad = bc$ por
$a$, obtém-se a fórmula de $d$.
\end{proof}

\begin{example}[Quarta proporcional]\label{ex:g7:prop:fourth}
Se $7$ livros iguais pesam $2.8$ kg, quanto pesam $12$ livros desses?
Tabela e regra da cruz:
\begin{center}
\renewcommand{\arraystretch}{1.3}
\begin{tabular}{l|cc}
livros & $7$ & $12$ \\
\hline
kg & $2.8$ & $m$ \\
\end{tabular}
\end{center}
\[
7 \times m = 2.8 \times 12
\quad\Longrightarrow\quad
m = \frac{2.8 \times 12}{7} = \frac{33.6}{7} = 4.8 \text{ kg}.
\]
Conferência pela unidade: um livro pesa $0.4$ kg, e doze pesam
$4.8$ kg.
\end{example}

\section{Porcentagens}

\begin{method}[Aplicar e encontrar porcentagens]\label{met:g7:prop:percent}
\begin{enumerate}
  \item \emph{Aplicar} $t\,\%$: multiplique por $\frac{t}{100}$
        ($12\,\%$ de $350$ é $0.12 \times 350 = 42$);
  \item \emph{encontrar} a porcentagem que uma parte representa:
        calcule $\frac{\text{parte}}{\text{todo}}$ e reescreva sobre
        $100$ ($21$ em $60$: $\frac{21}{60} = \frac{35}{100} =
        35\,\%$);
  \item pergunte sempre: \emph{porcentagem de quê?} O todo de
        referência importa.
\end{enumerate}
\end{method}

\begin{example}\label{ex:g7:prop:percent}
Numa escola, $180$ dos $400$ alunos são meninas. Porcentagem:
\[
\frac{180}{400} = \frac{180 \div 4}{400 \div 4} = \frac{45}{100}
= 45\,\% .
\]
No outro sentido: quantos são os $55\,\%$ de meninos?
$0.55 \times 400 = 220$. Confere: $180 + 220 = 400$.
\end{example}

\section{Escalas}

\begin{definition}[Escala]\label{def:g7:prop:scale}
Num mapa ou numa maquete de \emph{escala}\index{escala}
$\frac{1}{n}$ (escrita $1 : n$), as distâncias no mapa são
\omterm{def:g6:propdata:def}{proporcionais} às distâncias reais, com coeficiente $\frac1n$:
\[
\text{distância no mapa} = \frac{1}{n} \times \text{distância real},
\]
as duas medidas na \emph{mesma unidade}.
\end{definition}

\begin{example}\label{ex:g7:prop:scale}
Num mapa $1 : 50\,000$, duas cidades estão a $7$ cm uma da outra.
Distância real:
\[
7 \times 50\,000 = 350\,000 \text{ cm} = 3\,500 \text{ m} = 3.5
\text{ km}.
\]
Ao contrário, uma trilha de $12$ km mede no mapa: $12$ km $=
1\,200\,000$ cm, logo $1\,200\,000 \div 50\,000 = 24$ cm.
\end{example}

\begin{omfigure}
\begin{tikzpicture}
\begin{axis}[omaxis,
  width=8.4cm, height=5.6cm,
  xmin=0, xmax=8.6, ymin=0, ymax=34,
  xlabel={mapa (cm)}, ylabel={real (km)},
  xtick={2,4,6,8}, ytick={10,20,30}]
  \addplot[omDef, thick, domain=0:8.2] {4*x};
  \addplot[omDef, only marks, mark=*, mark size=1.5pt] coordinates
    {(2,8) (4,16) (6,24) (8,32)};
  \draw[dashed, omThm] (axis cs:7,0) -- (axis cs:7,28)
    -- (axis cs:0,28);
\end{axis}
\end{tikzpicture}

{\small Uma \omterm{def:g7:prop:scale}{escala} é uma \omterm{def:g7:prop:table}{proporcionalidade}: as distâncias no mapa
contra as distâncias reais se alinham numa reta que passa pela origem
(aqui, $1$ cm $\leftrightarrow$ $4$ km; o tracejado lê $7$ cm
$\leftrightarrow$ $28$ km).}
\end{omfigure}

\begin{remark}
Os gráficos são o teste mais rápido de \omterm{def:g7:prop:table}{proporcionalidade}
(\cref{ch:g6:propdata}): pontos numa reta \emph{que passa pela origem}
significam grandezas \omterm{def:g6:propdata:def}{proporcionais}; uma reta que não passa pela origem
(como a tarifa de um táxi com bandeirada fixa) significa que elas
\emph{não} são \omterm{def:g6:propdata:def}{proporcionais}, mesmo sendo uma reta. As funções afins,
no \cref{ch:g9:linfunc}, vão tornar isso preciso.
\end{remark}

\section{Exercícios}

\begin{exercise}[$\star$]\label{exo:g7:prop:1}
Que tabelas são \omterm{def:g7:prop:table}{tabelas de proporcionalidade}? Dê o coeficiente quando
ele existir.
\begin{center}
\renewcommand{\arraystretch}{1.2}
\begin{tabular}{l|ccc}
 & $3$ & $7$ & $11$ \\
\hline
 & $7.5$ & $17.5$ & $27.5$ \\
\end{tabular}
\qquad
\begin{tabular}{l|ccc}
 & $2$ & $5$ & $9$ \\
\hline
 & $6$ & $15$ & $28$ \\
\end{tabular}
\end{center}
\end{exercise}

\begin{exercise}[$\star$]\label{exo:g7:prop:2}
Complete usando a regra da cruz, escrevendo o cálculo: $5$ kg de batata
custam $6$ reais; quanto custam $8$ kg?
\end{exercise}

\begin{exercise}[$\star$]\label{exo:g7:prop:3}
Uma impressora imprime $36$ páginas em $3$ minutos. No mesmo ritmo,
quantas páginas em $5$ minutos? E quanto tempo para $96$ páginas?
\end{exercise}

\begin{exercise}[$\star$]\label{exo:g7:prop:4}
Calcule: $30\,\%$ de $250$; \quad $15\,\%$ de $60$; \quad o $t\,\%$
tal que $t\,\%$ de $80$ seja $20$.
\end{exercise}

\begin{exercise}[$\star$]\label{exo:g7:prop:5}
Em $25$ arremessos, uma jogadora de basquete acertou $18$. Qual é a
porcentagem de acerto dela? (Reescreva a \omterm{def:g6:fractions:def}{fração} sobre $100$.)
\end{exercise}

\begin{exercise}[$\star$]\label{exo:g7:prop:6}
Num mapa $1 : 25\,000$, uma trilha mede $9$ cm. Qual é o comprimento
real dela, em km? Um lago tem $2$ km de comprimento: que comprimento
ele tem no mapa?
\end{exercise}

\begin{exercise}[$\star$]\label{exo:g7:prop:7}
Um carrinho de coleção é feito na escala $1 : 43$. O carro real tem
$4.3$ m de comprimento. Que comprimento tem o carrinho, em cm?
\end{exercise}

\begin{exercise}[$\star\star$]\label{exo:g7:prop:8}
Uma receita usa $240$ g de chocolate para $8$ porções. Aline tem
$300$ g de chocolate. Para quantas porções isso dá (número inteiro de
porções)?
\end{exercise}

\begin{exercise}[$\star\star$]\label{exo:g7:prop:9}
O preço de uma jaqueta cai de $80$ para $60$ reais.
\begin{enumerate}
  \item De quanto é o desconto em reais? E em porcentagem do preço
        original?
  \item Mais tarde o preço volta de $60$ para $80$ reais. Explique por
        que essa alta \emph{não} é de $25\,\%$ --- calcule a
        porcentagem correta de aumento.
\end{enumerate}
\end{exercise}

\begin{exercise}[$\star\star$]\label{exo:g7:prop:10}
Duas linhas de uma tabela são \omterm{def:g6:propdata:def}{proporcionais}:
\begin{center}
\renewcommand{\arraystretch}{1.2}
\begin{tabular}{l|ccc}
 & $4$ & $x$ & $10$ \\
\hline
 & $y$ & $10.5$ & $17.5$ \\
\end{tabular}
\end{center}
Encontre o coeficiente e depois $x$ e $y$.
\end{exercise}

\begin{exercise}[$\star\star\star$]\label{exo:g7:prop:11}
Uma copiadora reduz os documentos a $80\,\%$ do tamanho deles. Um
\omterm{def:g6:lines:objects}{segmento} de $10$ cm é copiado e depois a \emph{cópia} é copiada de
novo com o mesmo ajuste.
\begin{enumerate}
  \item Que comprimento tem o \omterm{def:g6:lines:objects}{segmento} depois de uma cópia? E depois de
        duas?
  \item Por que a resposta depois de duas cópias não é $60\,\%$ do
        original? Que porcentagem única corresponde a duas cópias?
\end{enumerate}
\end{exercise}

\section{Problema: medir a Terra com uma vara}

\begin{problem}[{Problema de fim de semana --- as sombras são uma
proporcionalidade, e como Eratóstenes calculou o comprimento da Terra
em 240 a.C.}]
\label{pb:g7:prop:1}
Vinte e dois séculos atrás, sem telescópio, sem satélite e sem
calculadora, o bibliotecário de Alexandria mediu a Terra --- usando uma
vara, um poço, uma caravana de camelos e a matemática deste capítulo.
Este problema refaz os passos dele: primeiro a \omterm{def:g7:prop:table}{proporcionalidade} das
sombras, depois a famosa conta em si e, por fim, o que a resposta dele
parece em \omterm{def:g7:prop:scale}{escala} humana.

\medskip
\noindent\textbf{Parte I --- A sombra de uma vara.}
O Sol está tão longe que os raios dele chegam até nós paralelos uns aos
outros. Num dado momento, todos os objetos verticais e as sombras deles
são, portanto, \omterm{def:g6:propdata:def}{proporcionais}.
\begin{enumerate}
  \item Ao meio-dia, uma vara vertical de $1$ m projeta uma sombra de
        $0.4$ m, e a sombra de uma árvore mede $3.2$ m. Que altura tem
        a árvore (\cref{thm:g7:prop:cross})?
  \item No mesmo momento, uma torre projeta uma sombra de $26$ m: que
        altura tem a torre? E que comprimento tem a sombra de uma
        criança de $1.5$ m?
  \item Explique em uma frase por que todas essas contas só valem no
        \emph{mesmo momento} do dia.
  \item Conta a lenda que Tales mediu a Grande \omterm{def:g5:solids:def}{Pirâmide} esperando o
        momento em que \emph{a própria sombra dele ficasse exatamente
        do tamanho da altura dele}. Qual é o coeficiente de
        \omterm{def:g7:prop:table}{proporcionalidade} nesse momento, e que altura tem a \omterm{def:g5:solids:def}{pirâmide} se
        a ponta da sombra dela fica a $147$ m do ponto do chão que está
        logo abaixo do ápice?
  \item Resuma a Parte I: num dado momento, que grandeza é a mesma para
        a vara, a árvore, a torre e a \omterm{def:g5:solids:def}{pirâmide}
        (\cref{def:g7:prop:table})?
\end{enumerate}

\medskip
\noindent\textbf{Parte II --- A conta de Eratóstenes.}
Eratóstenes conhecia dois fatos. Na cidade de Siena, ao meio-dia do
solstício de verão, o Sol ficava \emph{exatamente a pino}: a luz
alcançava o fundo dos poços mais profundos, e as varas verticais não
projetavam sombra. Em Alexandria, $800$ km ao norte, no mesmo momento,
uma vara vertical \emph{projetava} sombra, e os \omterm{def:g3:shapes:shapes}{raios} do Sol faziam um
ângulo de $7.2^\circ$ com a vertical.
\begin{enumerate}[resume]
  \item Explique por que essas duas observações, juntas, provam que o
        chão de Siena e o chão de Alexandria não são paralelos --- isto
        é, que a superfície da Terra é curva. (O que raios paralelos
        fariam com duas varas numa Terra plana?)
  \item Num desenho da Terra redonda com raios de Sol paralelos, os
        $7.2^\circ$ de Alexandria reaparecem no centro da Terra, como o
        ângulo entre as direções das duas cidades. Faça o desenho e
        convença-se disso (a justificativa limpa usa os pares de
        ângulos iguais do \cref{ch:g7:triangles}).
  \item Que \omterm{def:g6:fractions:def}{fração} de uma volta completa são $7.2^\circ$?
  \item O arco de Siena até Alexandria ($800$ km) corresponde a
        $7.2^\circ$; o comprimento inteiro corresponde a $360^\circ$.
        Monte a \omterm{def:g7:prop:table}{proporcionalidade} e calcule o comprimento da Terra.
  \item Deduza o \omterm{def:g4:geometry:circle}{diâmetro} da Terra, usando $\pi \approx 3.14$
        (\cref{prop:g6:measure:circle}) e arredondando à centena de
        quilômetros mais próxima. (O valor moderno é $12\,742$ km ---
        quão perto chegou um homem com uma vara em 240 a.C.?)
\end{enumerate}

\medskip
\noindent\textbf{Parte III --- A Terra em \omterm{def:g7:prop:scale}{escala} humana.}
\begin{enumerate}[resume]
  \item Um globo é construído na escala $1 : 40\,000\,000$. Usando o
        comprimento encontrado na questão 9, mostre que o comprimento
        do globo é exatamente $1$ m. Qual é o \omterm{def:g4:geometry:circle}{diâmetro} dele, ao
        centímetro mais próximo?
  \item O Monte Branco sobe cerca de $4.8$ km. Converta $4.8$ km em
        centímetros, divida por $40\,000\,000$ e expresse a altura da
        montanha no globo em milímetros. O que a resposta diz sobre o
        quanto a Terra é de fato ``cheia de bolotas''?
  \item Um caminhante percorre $40$ km por dia. Nesse ritmo, quantos
        dias para o comprimento inteiro de Eratóstenes --- e isso dá,
        mais ou menos, quantos anos?
  \item A atmosfera respirável está concentrada, grosso modo, nos
        primeiros $10$ km acima do chão. Que porcentagem do \omterm{def:g3:shapes:shapes}{raio} da
        Terra (cerca de $6\,370$ km) é isso
        (\cref{met:g7:prop:percent})? Arredonde ao décimo de por
        cento.
  \item Os historiadores estimam que o valor anunciado por Eratóstenes,
        convertido para unidades modernas, pode ter sido de cerca de
        $42\,000$ km. Tomando $40\,000$ km como valor verdadeiro,
        calcule a porcentagem de erro dele. Conclua em uma frase sobre
        varas, poços e \omterm{def:g7:prop:table}{proporcionalidade}.
\end{enumerate}
\end{problem}
